\section{Tutorial 8: Gradient Flow and Instability}

\subsection{Objectives}

This tutorial develops intuition for gradient propagation in deep networks:

\begin{itemize}
    \item Understand how gradients evolve across layers
    \item Analyze vanishing and exploding gradients
    \item Study the role of activation functions
    \item Examine the effect of initialization
\end{itemize}

\medskip

\noindent
\textbf{Focus.}  
Understand gradient flow as a dynamical process across layers.

\subsection{Exercise 1: Gradient Through a Deep Linear Network}

Consider a deep linear network:
\[
f(x) = W_L W_{L-1} \cdots W_1 x.
\]

\medskip

\noindent
\textbf{Tasks.}
\begin{enumerate}
    \item Write the gradient $\frac{\partial L}{\partial x}$.
    \item Explain how depth affects gradient magnitude.
\end{enumerate}

\medskip

\noindent
\textbf{Solution.}

\[
\frac{\partial L}{\partial x}
= W_1^\top W_2^\top \cdots W_L^\top \frac{\partial L}{\partial f}.
\]

\medskip

\noindent
\begin{itemize}
    \item Gradient is a product of matrices
    \item Norm scales roughly as product of operator norms
\end{itemize}

\medskip

\noindent
\textbf{Insight.}  
Repeated multiplication leads to exponential growth or decay.

\subsection{Exercise 2: Scalar Case Intuition}

Consider:
\[
f(x) = a^L x.
\]

\medskip

\noindent
\textbf{Tasks.}
\begin{enumerate}
    \item Compute $\frac{df}{dx}$.
    \item Analyze behavior for:
    \begin{itemize}
        \item $|a| < 1$
        \item $|a| > 1$
    \end{itemize}
\end{enumerate}

\medskip

\noindent
\textbf{Solution.}

\[
\frac{df}{dx} = a^L.
\]

\begin{itemize}
    \item $|a| < 1$ $\Rightarrow$ vanishing gradients
    \item $|a| > 1$ $\Rightarrow$ exploding gradients
\end{itemize}

\medskip

\noindent
\textbf{Insight.}  
Depth amplifies small deviations multiplicatively.

\subsection{Exercise 3: Effect of Activation Functions}

Consider activations:

\begin{itemize}
    \item Sigmoid: $\sigma'(z) \leq 0.25$
    \item ReLU: $\sigma'(z) \in \{0,1\}$
\end{itemize}

\medskip

\noindent
\textbf{Tasks.}
\begin{enumerate}
    \item Why do sigmoid networks suffer from vanishing gradients?
    \item Why does ReLU mitigate this problem?
\end{enumerate}

\medskip

\noindent
\textbf{Solution.}

\begin{enumerate}
\item Sigmoid:
\begin{itemize}
    \item Derivatives are small ($<1$)
    \item Repeated multiplication $\Rightarrow$ vanishing gradients
\end{itemize}

\item ReLU:
\begin{itemize}
    \item Derivative is $1$ in active region
    \item Preserves gradient magnitude
\end{itemize}
\end{enumerate}

\medskip

\noindent
\textbf{Insight.}  
Activation choice directly affects gradient flow.

\subsection{Exercise 4: Variance Propagation (Conceptual)}

Assume:
\[
h_\ell = W_\ell h_{\ell-1},
\]
with random weights.

\medskip

\noindent
\textbf{Tasks.}
\begin{enumerate}
    \item What happens if weights are too large?
    \item What happens if weights are too small?
\end{enumerate}

\medskip

\noindent
\textbf{Solution.}

\begin{itemize}
    \item Large weights:
    \begin{itemize}
        \item Activations grow exponentially
        \item Exploding gradients
    \end{itemize}

    \item Small weights:
    \begin{itemize}
        \item Activations shrink
        \item Vanishing gradients
    \end{itemize}
\end{itemize}

\medskip

\noindent
\textbf{Insight.}  
Initialization controls signal propagation.

\subsection{Exercise 5: Gradient Flow as a Dynamical System}

\textbf{Question.}  
Why can we think of gradient propagation as a dynamical system?

\medskip

\noindent
\textbf{Discussion.}

\begin{itemize}
    \item Each layer transforms gradients
    \item Repeated transformations resemble iterative dynamics
\end{itemize}

\medskip

\noindent
\textbf{Insight.}  
Training stability depends on the stability of this dynamical system.

\subsection{Exercise 6: Practical Diagnosis}

\textbf{Tasks.}
\begin{enumerate}
    \item How can you detect vanishing gradients during training?
    \item How can you detect exploding gradients?
\end{enumerate}

\medskip

\noindent
\textbf{Solution.}

\begin{itemize}
    \item Vanishing:
    \begin{itemize}
        \item Near-zero gradients in early layers
        \item Slow learning
    \end{itemize}

    \item Exploding:
    \begin{itemize}
        \item Very large gradients
        \item Unstable loss or divergence
    \end{itemize}
\end{itemize}

\medskip

\noindent
\textbf{Insight.}  
Monitoring gradients is essential in practice.

\subsection{Exercise 7 (Discussion): Remedies}

\textbf{Question.}  
What techniques can mitigate gradient instability?

\medskip

\noindent
\textbf{Discussion points.}

\begin{itemize}
    \item Proper initialization
    \item ReLU-type activations
    \item Normalization methods
    \item Residual connections
\end{itemize}

\medskip

\noindent
\textbf{Insight.}  
Modern deep learning is largely about stabilizing gradient flow.

\subsection{Summary}

\begin{itemize}
    \item Gradient flow involves repeated transformations across layers
    \item Vanishing and exploding gradients arise from multiplicative effects
    \item Activation functions and initialization strongly influence stability
    \item Training deep networks requires controlling gradient dynamics
\end{itemize}

\medskip

\noindent
\textbf{Preparation for next lecture.}  
We now study activation functions and architectural design choices that improve gradient flow and expressivity.